%!TEX root = ../template.tex
%%%%%%%%%%%%%%%%%%%%%%%%%%%%%%%%%%%%%%%%%%%%%%%%%%%%%%%%%%%%%%%%%%%%
%% chapter3.tex
%% NOVA thesis document file
%%
%% Chapter with lots of dummy text
%%%%%%%%%%%%%%%%%%%%%%%%%%%%%%%%%%%%%%%%%%%%%%%%%%%%%%%%%%%%%%%%%%%%

\typeout{NT FILE chapter3.tex}%

\chapter{Project Planning}
\label{ch:project_planning}

\section{Tasks}
\label{sec:tasks}

The project can be split into tasks, described bellow.

\begin{enumerate}
    \item \textbf{Literature Review:} A comprehensive survey on the literature of the relevant topics is conducted, with the goal of building the theoretical foundations necessary for practical implementation. This literature review will address the topics of designing UAV controllers and reachability analysis, as well as any and all sub-topics branching out from these.
    \item \textbf{Definition of Requirements:} The problem will be explicitly formulated. The specific requirements of the "success" conditions of the study will be identified.
    \item \textbf{Open-Loop Design:} This phase consists of defining the mathematical model of both the drone and the missile dynamics. These models must be accurate, considering the full physical dynamics, and complete, considering aerodynamic effects. 
    \item \textbf{Defense-Intercept Algorithm Definition:} The reachability set-based control function is theoretically defined such that it meets the defense-intercept requirements previously laid out. The control function is implemented in the closed-loop of the drone model.
    \item \textbf{Controller Validation in Simulation:} The drone controller is systematically validated according to the MILS-SILS-PILS-HILS techniques. The simulation scenarios must also be defined such that they broach as wide as possible a range of real scenarios, by considering more or less sophisticated missile controllers and by considering adverse flight conditions.
    \item \textbf{Test Flight Validation:} The control function is loaded into the physical hardware of the drone and flight tests are conducted to validate its effectiveness and evaluate its performance. In these tests, the drone is facing off against a fixed-wing UAV that mimics a missile.
    \item \textbf{Writing Thesis:} The final document is written, describing the work done, evaluating the results and drawing up conclusions.
\end{enumerate}

\section{Planned Schedule}
\label{sec:planned_schedule}

Table~\ref{tab:gantt_chart} outlines the planned distribution of tasks over the time available, in the form of a Gantt Chart.

% Please add the following required packages to your document preamble:
% \usepackage{booktabs}
% \usepackage[table,xcdraw]{xcolor}
% Beamer presentation requires \usepackage{colortbl} instead of \usepackage[table,xcdraw]{xcolor} @{}p{1cm}p{1cm}p{1cm}p{1cm}p{1cm}p{1cm}p{1cm}p{1cm}p{1cm}p{1cm}p{1cm}@{}
\begin{table}[htbp]
\centering
\begin{tabular}{@{}p{1cm}p{1cm}p{1cm}p{1cm}p{1cm}p{1cm}p{1cm}p{1cm}p{1cm}p{1cm}p{1cm}@{}}
\toprule
\rowcolor[HTML]{C0C0C0} 
\multicolumn{2}{c|}{\cellcolor[HTML]{C0C0C0}2025} &
  \multicolumn{9}{c}{\cellcolor[HTML]{C0C0C0}2026} \\
\rowcolor[HTML]{EFEFEF} 
Nov &
  Dec &
  Jan &
  Feb &
  Mar &
  Apr &
  May &
  Jun &
  Jul &
  Aug &
  \multicolumn{1}{c}{\cellcolor[HTML]{EFEFEF}Sep} \\ \midrule
\rowcolor[HTML]{FFFFFF} 
\cellcolor[HTML]{34CDF9}T1 &
   &
   &
   &
   &
   &
   &
   &
   &
   &
   \\
\rowcolor[HTML]{FFFFFF} 
 &
  \cellcolor[HTML]{FFCC67}T2 &
   &
   &
   &
   &
   &
   &
   &
   &
   \\
\rowcolor[HTML]{FFFFFF} 
 &
   &
  \multicolumn{2}{c}{\cellcolor[HTML]{34CDF9}T3} &
   &
   &
   &
   &
   &
   &
   \\
\rowcolor[HTML]{FFFFFF} 
 &
   &
   &
   &
  \multicolumn{2}{c}{\cellcolor[HTML]{FFCC67}T4} &
   &
   &
   &
   &
   \\
\rowcolor[HTML]{FFFFFF} 
 &
   &
   &
   &
   &
   &
  \multicolumn{3}{c}{\cellcolor[HTML]{34CDF9}T5} &
   &
   \\
\rowcolor[HTML]{FFFFFF} 
 &
   &
   &
   &
   &
   &
   &
   &
   &
  \cellcolor[HTML]{FFCC67}T6 &
   \\
\rowcolor[HTML]{34CDF9} 
\cellcolor[HTML]{FFFFFF} &
  \cellcolor[HTML]{FFFFFF} &
  \cellcolor[HTML]{FFFFFF} &
  \cellcolor[HTML]{FFFFFF} &
  \cellcolor[HTML]{FFFFFF} &
  \multicolumn{6}{c}{\cellcolor[HTML]{34CDF9}T7} \\ \bottomrule
\end{tabular}
\caption{Gantt Chart}
\label{tab:gantt_chart}
\end{table}

\section{Prior Work}
\label{sec:prior_work}

In preparation for this study, the continuous mathematical model for the missile dynamics was derived, ignoring aerodynamical effects. This model is described in Subsection~\ref{subsec:missile}. This open-loop model was implemented in MATLAB, as well as the framework for simulations. The framework was made with a closed-loop structure in mind, running in discrete time with a sampling period of 0.1 seconds, so that control inputs may be changed time-step to time-step. Despite this the model used in the simulation is continuous, as the MATLAB function ode45 is used to compute the values in the following time-step, and the sampling period is sufficiently small to ignore distortions. \\

Figure~\ref{fig:prev_work} shows a missile simulation meant to verify if the implemented dynamics behave as expected. The initial position of the base of the missile is at the origin, velocity is zero, angular velocity is zero and the rotation matrix is identity, so that the missile sits upright. The blue vector represents the missile body, magenta are the thrust acceleration vectors acting on the missile, and red is the trajectory of the missile until this instant.

\begin{figure}[htbp]
  \centering
  \includegraphics[width=0.8\linewidth]{ch3_prev_work}%
  \caption{Simulation of previous work.}
  \label{fig:prev_work}
\end{figure}